% GENERATED FILE. Edit canonical lessons, then run npm run generate:books.
% canonical-source-hash: fnv1a64:3b3d26ee7175fe52
% canonical-lessons: LA-C533-familiaris, LA-C533-familiaritas, LA-C533-necessitudo, LA-C533-diligo, LA-C533-amabilis

\chapter{Close friends}
\label{ch:la-close-friends}
\hlchaptermodality{533}

\begin{chapteropening}
\textbf{By the end of this chapter:} \emph{I can talk about close friends and the ties between people.}

Complete the last lesson of chapter 533: I can talk about close friends and the ties between people.
\end{chapteropening}

\section[familiāris, familiāris]{familiāris, familiāris — a close friend}

\label{lesson:LA-C533-familiaris}



\begin{quote}
Before the new one: say the Latin for free-born, then the Latin for a felt cap.
\end{quote}

\subsection*{You'll want to know: familiāris, familiāris}
\textbf{familiāris, familiāris} — \textquotedblleft{}a close friend\textquotedblright{}.

\begin{quote}
Hic vir familiāris patris meī est.
\end{quote}

\begin{quote}
\emph{This man is a close friend of my father's.}
\end{quote}

As an adjective (\emph{familiāris, familiāre}), of the household; as a noun, an intimate friend. English \emph{familiar}.

\begin{cousinweb}
\textbf{familiāris} is built on \emph{familia} \textquotedblleft{}household\textquotedblright{} (Lewis \& Short; EtymDB). English \emph{familiar} comes from it.
\end{cousinweb}

\begin{grammarlens}[title={What you've built}]
One more word for telling, explaining and describing at B1.
\end{grammarlens}

\subsection*{Your turn}
\begin{itemize}
\raggedright
  \item \emph{Say it:} \emph{familiāris, familiāris}
  \item \emph{Say it:} \emph{familiāris, familiāris}, once more
  \item \emph{Say it:} a sentence of your own with \emph{familiāris, familiāris}
  \item \emph{Recall:} say the Latin for free-born, then the Latin for a felt cap, then say \emph{familiāris, familiāris} again
\end{itemize}

\begin{tcolorbox}[breakable,colback=teal!4,colframe=teal!35!black,title={Before you move on}]
What does familiāris, familiāris mean? (\textquotedblleft{}A close friend\textquotedblright{}.) Say it once more.
\end{tcolorbox}
\section[familiāritās, familiāritātis]{familiāritās, familiāritātis — close friendship, intimacy}

\label{lesson:LA-C533-familiaritas}



\begin{quote}
Before the new one: say the Latin for a felt cap, then the Latin for a close friend.
\end{quote}

\subsection*{You'll want to know: familiāritās, familiāritātis}
\textbf{familiāritās, familiāritātis} — \textquotedblleft{}close friendship, intimacy\textquotedblright{}.

\begin{quote}
Inter eōs magna familiāritās est.
\end{quote}

\begin{quote}
\emph{There is a close friendship between them.}
\end{quote}

Feminine, third declension. English \emph{familiarity}.

\begin{cousinweb}
\textbf{familiāritās} is built on \emph{familiāris} \textquotedblleft{}intimate, of the household\textquotedblright{} (Lewis \& Short). English \emph{familiarity} comes from it.
\end{cousinweb}

\begin{grammarlens}[title={What you've built}]
One more word for telling, explaining and describing at B1.
\end{grammarlens}

\subsection*{Your turn}
\begin{itemize}
\raggedright
  \item \emph{Say it:} \emph{familiāritās, familiāritātis}
  \item \emph{Say it:} \emph{familiāritās, familiāritātis}, once more
  \item \emph{Say it:} a sentence of your own with \emph{familiāritās, familiāritātis}
  \item \emph{Recall:} say the Latin for a felt cap, then the Latin for a close friend, then say \emph{familiāritās, familiāritātis} again
\end{itemize}

\begin{tcolorbox}[breakable,colback=teal!4,colframe=teal!35!black,title={Before you move on}]
What does familiāritās, familiāritātis mean? (\textquotedblleft{}Close friendship, intimacy\textquotedblright{}.) Say it once more.
\end{tcolorbox}
\section[necessitūdō, necessitūdinis]{necessitūdō, necessitūdinis — a close bond, tie}

\label{lesson:LA-C533-necessitudo}



\begin{quote}
Before the new one: say the Latin for a close friend, then the Latin for close friendship.
\end{quote}

\subsection*{You'll want to know: necessitūdō, necessitūdinis}
\textbf{necessitūdō, necessitūdinis} — \textquotedblleft{}a close bond, tie\textquotedblright{}.

\begin{quote}
Magna necessitūdō mihi cum frātre tuō est.
\end{quote}

\begin{quote}
\emph{I have a close bond with your brother.}
\end{quote}

Feminine, third declension, from \emph{necesse}. Also need, but in letters it often means a tie of family or friendship.

\begin{cousinweb}
\textbf{necessitūdō} is built on \emph{necesse} \textquotedblleft{}necessary\textquotedblright{} (Lewis \& Short): first need, then the ties that bind people.
\end{cousinweb}

\begin{grammarlens}[title={What you've built}]
One more word for telling, explaining and describing at B1.
\end{grammarlens}

\subsection*{Your turn}
\begin{itemize}
\raggedright
  \item \emph{Say it:} \emph{necessitūdō, necessitūdinis}
  \item \emph{Say it:} \emph{necessitūdō, necessitūdinis}, once more
  \item \emph{Say it:} a sentence of your own with \emph{necessitūdō, necessitūdinis}
  \item \emph{Recall:} say the Latin for a close friend, then the Latin for close friendship, then say \emph{necessitūdō, necessitūdinis} again
\end{itemize}

\begin{tcolorbox}[breakable,colback=teal!4,colframe=teal!35!black,title={Before you move on}]
What does necessitūdō, necessitūdinis mean? (\textquotedblleft{}A close bond, tie\textquotedblright{}.) Say it once more.
\end{tcolorbox}
\section[dīligō, dīligere]{dīligō, dīligere — to hold dear, esteem}

\label{lesson:LA-C533-diligo}



\begin{quote}
Before the new one: say the Latin for close friendship, then the Latin for a close bond.
\end{quote}

\subsection*{You'll want to know: dīligō, dīligere}
\textbf{dīligō, dīligere} — \textquotedblleft{}to hold dear, esteem\textquotedblright{}.

\begin{quote}
Parentēs līberōs suōs dīligunt.
\end{quote}

\begin{quote}
\emph{Parents hold their children dear.}
\end{quote}

\emph{dīligō, dīligere, dīlēxī, dīlēctum}. Love by choice and esteem, calmer than \emph{amō}. \emph{Dīligēns}, already known, comes from it.

\begin{cousinweb}
\textbf{dīligō} is built on \emph{dis-} \textquotedblleft{}apart\textquotedblright{} and \emph{legō} \textquotedblleft{}to choose\textquotedblright{} (Lewis \& Short; EtymDB): to single out. English \emph{diligent} and \emph{predilection} come from the family.
\end{cousinweb}

\begin{grammarlens}[title={What you've built}]
One more word for telling, explaining and describing at B1.
\end{grammarlens}

\subsection*{Your turn}
\begin{itemize}
\raggedright
  \item \emph{Say it:} \emph{dīligō, dīligere}
  \item \emph{Say it:} \emph{dīligō, dīligere}, once more
  \item \emph{Say it:} a sentence of your own with \emph{dīligō, dīligere}
  \item \emph{Recall:} say the Latin for close friendship, then the Latin for a close bond, then say \emph{dīligō, dīligere} again
\end{itemize}

\begin{tcolorbox}[breakable,colback=teal!4,colframe=teal!35!black,title={Before you move on}]
What does dīligō, dīligere mean? (\textquotedblleft{}To hold dear, esteem\textquotedblright{}.) Say it once more.
\end{tcolorbox}
\section[amābilis, amābile]{amābilis, amābile — lovable, delightful}

\label{lesson:LA-C533-amabilis}



\begin{quote}
Before the new one: say the Latin for a close bond, then the Latin for to hold dear.
\end{quote}

\subsection*{You'll want to know: amābilis, amābile}
\textbf{amābilis, amābile} — \textquotedblleft{}lovable, delightful\textquotedblright{}.

\begin{quote}
Fīliola tua amābilis est.
\end{quote}

\begin{quote}
\emph{Your little daughter is lovable.}
\end{quote}

Two endings, from \emph{amō} with \emph{-bilis}.

\begin{cousinweb}
\textbf{amābilis} is built on \emph{amō} \textquotedblleft{}to love\textquotedblright{} with \emph{-bilis} \textquotedblleft{}able to be\textquotedblright{} (Lewis \& Short; EtymDB).
\end{cousinweb}

\begin{grammarlens}[title={What you've built}]
One more word for telling, explaining and describing at B1.
\end{grammarlens}

\subsection*{Your turn}
\begin{itemize}
\raggedright
  \item \emph{Say it:} \emph{amābilis, amābile}
  \item \emph{Say it:} \emph{amābilis, amābile}, once more
  \item \emph{Say it:} a sentence of your own with \emph{amābilis, amābile}
  \item \emph{Recall:} say the Latin for a close bond, then the Latin for to hold dear, then say \emph{amābilis, amābile} again
\end{itemize}

\begin{tcolorbox}[breakable,colback=teal!4,colframe=teal!35!black,title={Before you move on}]
What does amābilis, amābile mean? (\textquotedblleft{}Lovable, delightful\textquotedblright{}.) Say it once more.
\end{tcolorbox}
